\documentclass[11pt]{article}
\usepackage[T1]{fontenc}
\usepackage{lmodern,amsmath,amssymb,microtype}
\usepackage[margin=1in]{geometry}
\usepackage[colorlinks=true,urlcolor=blue]{hyperref}
\hypersetup{pdfauthor={Wentao Li},pdftitle={The failure probability of granular synchrony on random networks}}
\title{The failure probability of granular synchrony on random networks}
\author{Wentao Li\\[2pt]\small Prepared with assistance from OpenAI Codex}
\date{November 2026}
\begin{document}
\maketitle
\begin{abstract}
On a network whose synchronous links appear independently, we compute the probability of violating the graph condition for granular synchrony to relative error $o(1)$. At linear honest-set size and minimum link probability $p_*$ satisfying $np_*/\log n\to\infty$, almost every failing topology has a unique separating partition, up to exchange of its sides. The formula allows nonidentical link probabilities. For a fixed fault fraction in the homogeneous model it simplifies to one multinomial term. Below a specified density the expected witness count can instead overestimate the failure probability by an unbounded factor.
\end{abstract}

\section{What constitutes a failure}
Fix $n/3\le f<n/2$, let $h=n-2f$, and let $S$ be the undirected graph of synchronous links on $n$ parties. Other links have the slower timing allowed by the protocol. The graph condition of Giridharan et al. [1] requires that, after deleting any $f$ parties, every set of $h$ remaining parties reaches at least $f+1$ parties along synchronous links. Appan, Heath and Ren [2] identify the same graph condition for secure computation with guaranteed output delivery in their model. We randomize only the assignment of synchronous links before applying the worst-case condition over fault sets.

A \emph{witness} is an unordered pair of disjoint shores $\{X,Y\}$, each of size at least $h$, whose union has size $n-f$, with no synchronous edge between them. The leftover $f$ vertices form $F$. To see the equivalence, start with a set of $h$ honest vertices reaching at most $f$ after deleting $F$. Its entire reachable set is one shore, leaving at least $n-f-f=h$ vertices on the other. Conversely, any $h$ vertices of one shore reach at most $n-f-h=f$ vertices. This separator description follows from the published condition. Let $Z$ count witnesses.

\section{Relative probability and uniqueness}
Include each synchronous edge $uv$ independently with probability $p_{uv}$. Write $p_*=\min_{uv}p_{uv}$ and let $\mathcal P$ be the unordered shore pairs satisfying the size requirements, without an edge requirement. Set
\[
 w(X,Y)=\prod_{u\in X,v\in Y}(1-p_{uv}),\qquad
 \mu=\sum_{\{X,Y\}\in\mathcal P}w(X,Y)=\mathbb E Z.
\]

\paragraph{Theorem.} Suppose $h\ge\epsilon n$ for a fixed $\epsilon>0$ and $np_*/\log n\to\infty$. Uniformly under these bounds,
\[
        (1-o(1))\mu\ \le\ \Pr(Z=1)\ \le\
        \Pr(Z\ge1)\ \le\ \mu .
\]
If $\mu>0$, a failing topology thus has one witness with conditional probability $1-o(1)$. The randomness describes networks; it is independent of any randomness in the protocols.

\paragraph{Proof.}
Fix a candidate $\{X,Y\}$ of positive weight and condition on it being a witness. All edges outside $X\times Y$ retain their original independent laws. Write $q_*=1-p_*$. Consider another candidate with arbitrarily ordered shores $A,B$. In each shore choose a largest intersection with the three original classes $X,Y,F$. Each chosen intersection has size at least $h/3$.

If the two chosen classes are anything other than $X,Y$ in some order, every edge between the two intersections is unconditioned. This includes the possibility that both intersections lie in the same class, since $A$ and $B$ are disjoint. At least $h^2/9$ edges must be absent. A union bound over at most $3^n$ ordered assignments gives a contribution at most $3^n q_*^{h^2/9}$.

Otherwise orient $A,B$ to have $|A\cap X|,|B\cap Y|\ge h/3$. Write $r=|A\setminus X|+|B\setminus Y|$. If $r=0$, equality of the total sizes forces $A=X,B=Y$. For $r>0$, edges from $A\setminus X$ to $B\cap Y$, and from $B\setminus Y$ to $A\cap X$, are unconditioned and form disjoint sets. At least $rh/3$ such edges must be absent. At most $2r$ vertices have changed class: the $r$ entries into shores from a different class account for all changes except moves into $F$, and there are at most $r$ of those by preservation of $|F|$. At most $(2n+1)^{2r}$ assignments can make these changes, by listing changed vertices and their new classes, with a padding symbol. Hence the conditional probability of any second witness is at most
\[
   3^n q_*^{h^2/9}
   +2\sum_{r=1}^{n}\big((2n+1)^2q_*^{h/3}\big)^r=o(1).
\]
The last estimate uses $q_*\le e^{-p_*}$ and is uniform in the planted shores and individual probabilities. Summing $w(X,Y)$ times the conditional probability of uniqueness counts each one-witness topology exactly once. The upper bound follows from $1_{\{Z\ge1\}}\le Z$. \hfill$\square$

\section{Identical link probabilities}
Write $q=1-p$, with $0<p<1$. The exact first moment appearing in the theorem is
\[
  \mu=\tfrac12\binom nf
       \sum_{a=h}^{f}\binom{n-f}{a}\,q^{a(n-f-a)} .
\]
The half removes the two possible shore orders, including for equal shore sizes. If $f/n\to\rho\in(1/3,1/2)$ and $np/\log n\to\infty$, this gives
\[
  \Pr(Z\ge1)\sim\frac{n!}{h!\,(f!)^2}\,q^{hf}.
\]
Indeed, the term with smaller shore $h+i$ relative to that with smaller shore $h$ is at most
$(f/h)^i q^{i(f-h-i)}$, for $1\le i\le(f-h)/2$. Its sum is $o(1)$ since $f-h$ is linear in $n$ and $np\to\infty$. At $n=3f$ the answer is instead asymptotic to $n!\,q^{f^2}/(2(f!)^3)$, since there is just one possible pair of shore sizes. Near this boundary the full sum applies.

\paragraph{When the first moment overcounts.}
Keep either $n=3f$ or $f/n\to\rho\in(1/3,1/2)$, and suppose $np\to\infty$ but $f q^h\to\infty$. Then $\Pr(Z\ge1)/\mu\to0$. To see this, sample a candidate with probability $w(X,Y)/\mu$, then sample the graph conditional on that witness. Under this experiment $\mathbb E[1/Z]=\Pr(Z\ge1)/\mu$, by summing over witnesses. The preceding ratio estimate uses only $np\to\infty$, so the sampled witness has sizes $h,f$ with probability $1-o(1)$. Call the smaller shore $X$ (choose either in a tie). Each of the $f$ vertices in $F$ has no neighbor in $X$ with probability $q^h$, independently across those vertices. With probability $1-o(1)$ some such vertex exists. Replacing in turn each of the $f$ vertices on the other shore by that vertex gives $f$ further distinct witnesses. Hence $1/Z\to0$ in probability in this experiment; as $1/Z\le1$, its expectation tends to zero. This gives a sparse regime where simply counting obstructions gives the wrong relative probability.

\section*{Sources and status}
[1] N. Giridharan, I. Abraham, N. Crooks, K. Nayak and L. Ren, \href{https://doi.org/10.4230/LIPIcs.DISC.2024.30}{Granular Synchrony}, DISC 2024, Theorem 13.\par
[2] A. Appan, D. Heath and L. Ren, \href{https://eprint.iacr.org/2025/1327}{Randomized Agreement, Verifiable Secret Sharing and Multi-Party Computation in Granular Synchrony}, ASIACRYPT 2025.\par
[3] D. Karger, \href{https://people.csail.mit.edu/karger/Papers/lastrel-stoc.pdf}{A Phase Transition and a Quadratic Time Unbiased Estimator for Network Reliability}, STOC 2020. This relates failure probability and expected cut counts for ordinary all-terminal disconnection. Here each witness also designates $f$ excluded vertices and two lower-bounded shores.
\par\medskip
We formulated the random-topology question. Independent review of the proof and priority of this application are pending.
\end{document}
